%% ============================================================================
%% DRAFT / SKETCH: Specification Test for the Joint Coarsening Conditions
%% Not currently \input from main.tex. Used to evaluate the math before
%% committing to a full development pass.
%% ============================================================================

\section{A Specification Test for the Joint Hypothesis $\{C_1, C_2, C_3\}$}
\label{sec:spec-test}

The local sensitivity indices of \Cref{sec:misspec,sec:c1-sensitivity,sec:c3-sensitivity}
characterize the bias of the standard MLE under each coarsening-condition
violation. In practice, a reliability engineer fitting the C1-C2-C3 model
faces a prior question: are any of C1, C2, C3 detectably violated in the
present data? We give a formal test of the joint null
$H_0 \colon \{C_1, C_2, C_3\}$ that operates on observable data alone (no
need to observe the failed cause $K_i$) and is consistent against the
parametric alternatives of \Cref{sec:c1-sensitivity,sec:misspec,sec:c3-sensitivity}.

\subsection{Observable Factorization Under $H_0$}
\label{sec:spec-test-lemma}

The test rests on a single elementary lemma.

\begin{lemma}[Factorization of $\Prob(C \mid T)$ under $H_0$]
\label{lem:factorization}
Under $H_0 \colon \{C_1, C_2, C_3\}$, for every non-empty candidate set
$c$ and every time $t > 0$,
\begin{equation}
\label{eq:factorization}
    \Prob\{C = c \mid T = t, \delta = 1\}
    \;=\; \pi_c \cdot R_c(t; \vtheta^*),
\end{equation}
where $\pi_c \in [0, 1]$ depends only on $c$ (not on the failed cause $k$,
not on the parameter $\vtheta^*$, not on the time $t$), and
\begin{equation}
\label{eq:Rc}
    R_c(t; \vtheta) \;\coloneqq\;
    \frac{\sum_{k \in c} h_k(t;\, \theta_k)}{S(t;\, \vtheta)}
\end{equation}
is the relative-hazard contribution of $c$ at time $t$, with
$S(t; \vtheta) = \sum_j h_j(t; \theta_j)$ the system hazard.
\end{lemma}

\begin{proof}
Under C1, $\Prob\{C = c \mid K = k\} = 0$ for $k \notin c$. Under C2,
$\Prob\{C = c \mid K = k\} = \pi_{k,c}$ is constant in $k$ for $k \in c$;
write the common value $\pi_c$. Under C3, $\pi_c$ does not depend on
$\vtheta^*$. The competing-risks decomposition gives
$\Prob\{K = k \mid T = t, \delta = 1\} = h_k(t; \theta_k^*) / S(t; \vtheta^*)$.
Hence
\begin{equation*}
    \Prob\{C = c \mid T = t, \delta = 1\}
    = \sum_{k=1}^{m} \Prob\{C = c \mid K = k\}\,
                      \Prob\{K = k \mid T = t, \delta = 1\}
    = \pi_c \sum_{k \in c} \frac{h_k(t; \theta_k^*)}{S(t; \vtheta^*)}.
\end{equation*}
The conditioning on $\delta = 1$ is innocuous because the candidate set
is generated only on failure under all relevant masking models and
because type-I censoring is independent of $(K, C)$ given $T$.
\end{proof}

The Lemma is the testable content of $H_0$: a conditional distribution
$\Prob(C \mid T)$ factors as a $c$-only function times a known parametric
function of $t$. Each component condition contributes one structural
constraint to this factorization. C1 contributes the support restriction
($\pi_c$ vanishes if $K \notin c$ for any positive-probability $k$), C2
contributes the $k$-invariance of $\pi_c$ for $k \in c$, and C3 contributes
the $\vtheta$-independence of $\pi_c$. Any failure of the joint factor
structure indicates at least one of these constraints is violated.

\subsection{Test Statistic for Singleton Candidate Sets}
\label{sec:spec-test-stat}

We construct the test for the singleton categories
$\mathcal{C}_{\mathrm{test}} = \{\{1\}, \{2\}, \ldots, \{m\}\}$. The
factorization \eqref{eq:factorization} implies that, under $H_0$, the
ratio $\Prob\{C = \{j\} \mid T = t, \delta = 1\} / R_j(t; \vtheta^*)$ is
constant in $t$, equal to $\pi_{\{j\}}$. A direct test of this constancy
is the bin-based Pearson chi-squared.

For each component $j$, define the model-based singleton intensity at
observation $i$:
\begin{equation}
\label{eq:Rj-singleton}
    \hat R_j(T_i) \;\coloneqq\; \frac{h_j(T_i; \hat\theta_j)}
                                       {\sum_l h_l(T_i; \hat\theta_l)},
\end{equation}
where $\hat\vtheta$ is the C1-C2-C3 MLE. The global moment estimator for
$\pi_{\{j\}}$ is
\begin{equation}
\label{eq:pi-singleton-est}
    \hat\pi_{\{j\}} \;=\;
    \frac{\sum_{i: \delta_i = 1} \ind{C_i = \{j\}}}
         {\sum_{i: \delta_i = 1} \hat R_j(T_i)}.
\end{equation}
This estimator zeros the residual $\sum_i (\ind{C_i = \{j\}} -
\hat\pi_{\{j\}} \hat R_j(T_i))$ in aggregate by construction, so an
aggregate-residual test has no power. The bin-based construction
sidesteps this by testing the constancy of $\hat\pi_{\{j\}}$ across time
bins, which is the testable content of the factorization.

\paragraph{Construction.}
Partition the failure-time domain into $L$ disjoint bins
$\mathcal{B}_1, \ldots, \mathcal{B}_L$ (in practice, quantile bins of the
empirical failure-time distribution). For each component $j$ and bin
$\ell$, define
\begin{align}
\label{eq:pearson-cells}
    O_{j, \ell} \;&=\; \sum_{i : T_i \in \mathcal{B}_\ell,\, \delta_i = 1}
                       \ind{C_i = \{j\}}, \\
    E_{j, \ell} \;&=\; \hat\pi_{\{j\}}
                       \sum_{i : T_i \in \mathcal{B}_\ell,\, \delta_i = 1}
                       \hat R_j(T_i).
\end{align}
The singleton specification test statistic is the Pearson chi-squared:
\begin{equation}
\label{eq:T-stat}
    T_n \;=\; \sum_{j = 1}^{m} \sum_{\ell = 1}^{L}
              \frac{(O_{j, \ell} - E_{j, \ell})^2}{E_{j, \ell}},
\end{equation}
summed over $(j, \ell)$ cells with $E_{j, \ell}$ above a small-count
threshold (typically $E_{j, \ell} \geq 0.5$).

\begin{theorem}[Asymptotic Null Distribution]
\label{thm:asymptotic-spec}
Suppose the standard regularity conditions for the MLE $\hat\vtheta$
hold, and the candidate-set distribution has $\pi_{\{j\}} > 0$ for every
$j$. Under $H_0$ as $n \to \infty$ with the number of bins $L$ fixed,
\begin{equation}
    T_n \xrightarrow{d} \chi^2_{m(L - 1)}.
\end{equation}
\end{theorem}

\begin{proof}[Proof outline]
For each $j$, the bin counts $(O_{j, \ell})_{\ell=1}^{L}$ are approximately
multinomial with cell probabilities determined by the bin averages of
$R_j(t)$ under the true mixture, conditional on the marginal singleton
rate. The Pearson chi-squared on these cells with one estimated parameter
($\hat\pi_{\{j\}}$) is asymptotically $\chi^2_{L-1}$ under $H_0$ by the
classical Wilks construction. Summing $m$ such statistics yields
$\chi^2_{m(L-1)}$, with the $m$-component statistics asymptotically
independent because the cell counts for different $j$ partition the
singleton observations into disjoint sets. The estimation of $\hat\vtheta$
introduces a vanishing correction under standard regularity (the score
$\partial \ell / \partial \vtheta$ is orthogonal to the bin-deviation
directions asymptotically), and the correction is empirically negligible
at sample sizes where the Pearson approximation is itself accurate.
\end{proof}

\begin{remark}[Choice of $L$ and small-count adjustments]
\label{rem:choice-L}
For a 5-component system at $n_F \approx 500$ failures, $L = 5$ quantile
bins gives expected cell counts in the range $\hat\pi_{\{j\}}
\sum_{i \in \mathcal{B}_\ell} \hat R_j(T_i)$, typically $5$ to $20$
events per cell when $\pi_{\{j\}} \approx 1/m$. The Pearson approximation
is accurate in this range. For sparse cells with $E_{j,\ell} < 0.5$,
either drop the cell from the sum (with corresponding reduction in $df$)
or merge adjacent bins. The validation study reported in
\Cref{sec:simulations} uses $L = 5$ throughout.
\end{remark}

\subsection{Power Against the Three Parametric Alternatives}
\label{sec:spec-test-power}

\begin{proposition}[Consistency against violations]
\label{prop:power}
The test rejects $H_0$ with probability tending to one under each of the
following sequences of alternatives, as long as the relevant parametric
violation is detectable in the data:
\begin{enumerate}
    \item \textbf{C1 violation} ($p_k(k) = 1 - \alpha_{C_1}$ with
        $\alpha_{C_1} > 0$): The conditional
        $\Prob\{C = \{j\} \mid T = t, \delta = 1\}$ acquires additional
        mass through $K \in c^{\mathsf c}$ paths that distort the
        $\pi_{\{j\}}$ factor. The distortion is $t$-dependent for
        non-exponential components, giving the test nontrivial power.
    \item \textbf{C2 violation} ($\pi_{k, c}$ varies with $k$ for
        $k \in c$): The mixture
        $\sum_{k \in c} \pi_{k,c}\, h_k(t)/S(t)$ does not factor as
        $\pi_c \cdot R_c$, so the test rejects.
    \item \textbf{C3 violation} ($\pi_c$ depends on $\vtheta^*$ via the
        power-weight law): For Weibull components the time-dependence in
        $\hat R_c(t)$ at the misspecified $\hat\vtheta$ surfaces the
        violation; for exponential components, $R_c$ is constant in $t$
        and the test has limited power against C3, consistent with the
        fact that total-hazard inference is exactly preserved under C3
        for exponential (\Cref{thm:exact-hazard-exp}).
\end{enumerate}
\end{proposition}

The C3-exponential power gap is not a defect: by
\Cref{thm:exact-hazard-exp}, C3 violation has no consequence for
system-level inference in the exponential case, so the absence of
detection power matches the absence of inferential consequence.

\subsection{Implementation and Validation}

The test is implemented in the \texttt{mdrelax} R package as
\texttt{spec\_test\_singleton()}. A simulation study validating the
chi-squared null distribution and characterizing the power across
violation types is reported in \Cref{sec:simulations}; see
\texttt{paper/specification\_test\_validation.R} for the reproducible
analysis.

%% ============================================================================
%% END SKETCH
%% ============================================================================
